% MACH-2: the same four-socket node, but the HARP ASICs extend cache coherence
% beyond the node, so many nodes together present one huge shared address
% space. The arrows leaving the top and bottom go to the other nodes.
\documentclass[dvisvgm,border=3pt]{standalone}
\input{preamble.tex}
\input{numa.tex}
\begin{document}
\begin{tikzpicture}[x=1cm,y=1cm]

  \numanodebase

  \node[container,minimum width=14mm,minimum height=7mm] (ht) at (4.32,4.80) {HARP};
  \node[container,minimum width=14mm,minimum height=7mm] (hb) at (4.32,0.38) {HARP};

  % anchored on the corners so the links read as diagonals, as in the original
  \draw[link] (ht.south west) -- (c0.north east);
  \draw[link] (ht.south east) -- (c1.north west);
  \draw[link] (hb.north west) -- (c2.south east);
  \draw[link] (hb.north east) -- (c3.south west);

  % NUMAlink carries coherence traffic on to the neighbouring nodes
  \draw[link,draw=blu,line width=1.3pt] (4.32,5.20) -- (4.32,5.85);
  \draw[link,draw=blu,line width=1.3pt] (4.32,0.00) -- (4.32,-0.65);

  \node[text=fg,font=\sffamily\small] at (4.3,-1.25) {SGI NUMAlink architecture};

\end{tikzpicture}
\end{document}
